% frontmatter/mapa-del-curso.tex -- el año entero en una tabla (el plan
% completo, semana a semana, está en notes/01-plan.md).
\section*{Mapa del curso}

\begin{center}
\renewcommand{\arraystretch}{1.4}
\begin{tabularx}{\linewidth}{@{}c c c c X@{}}
\toprule
\textbf{Trimestre} & \textbf{Días} & \textbf{Época} & \textbf{Números} & \textbf{Qué pasa} \\
\midrule
1 & 1--65 & otoño & 0 al 10 & Un número nuevo cada semana: se traza, se
cuenta, se colorea y se busca. El cumpleaños de Lucía tiene siete velas,
y la Navidad cierra el trimestre. \\
2 & 66--130 & invierno & hasta el 20 & La decena, y los números del 11
al 20; más, menos e igual; el anterior y el siguiente; las primeras
sumas, juntando dibujos. \\
3 & 131--195 & primavera & hasta el 50 & Sumar y restar hasta 10;
partir un número en dos; contar de 2 en 2, de 5 en 5 y de 10 en 10. \\
4 & 196--260 & verano & hasta el 100 & Las decenas hasta el 100; sumas y
restas hasta 20; el dinero, la hora y medir, en el pueblo de la abuela. \\
\bottomrule
\end{tabularx}
\end{center}

\vspace{6mm}
Cada trimestre tiene 13 semanas de 5 días -- también durante las
vacaciones: el cuaderno sigue siendo una página por día de lunes a
viernes en Navidad y en verano, como \emph{Aprendo a leer}. Un número
nuevo no llega hasta que el anterior ya se ha trazado, contado,
coloreado y buscado durante toda una semana.
\newpage
